\section{Introduction}
\label{sec:introduction}

Among the most widely discussed characteristics of text generated by large language models is the frequency of the em dash. Since the broad deployment of instruction-tuned models beginning in late 2022, users across writing communities, technology forums, and social media have identified em dash frequency as a marker of machine-generated prose---a ``tell'' that has become the subject of extensive online discussion, mockery, and ad hoc detection heuristics.

Separately, a parallel conversation has developed in developer and AI research communities: the observation that LLMs ``think in markdown.'' Models default to structured output---headers, bullet points, bold emphasis, numbered lists---even when such formatting is not requested. This tendency is widely attributed to the prevalence of markdown-formatted text in training corpora and is discussed primarily as a user experience or prompt engineering concern.

These two conversations describe the same phenomenon at different granularities, but they have not been connected. The writing communities observing em dash frequency do not frame their complaint in terms of training data composition or structural formatting. The developer communities analyzing markdown behavior do not extend their analysis to prose-level artifacts like punctuation patterns. A gap exists between the aesthetic observation (``AI overuses em dashes'') and the technical observation (``AI thinks in markdown''), and this paper fills it.

We propose that the em dash is \textbf{markdown leaking into prose}---the smallest surviving unit of the structural orientation that LLMs acquire from markdown-saturated training corpora. When models are constrained to produce natural prose rather than formatted output, the structural impulse is largely suppressed: headers, bullets, and bold emphasis are withheld. But the em dash survives in some models because it occupies a unique dual-register position: it is simultaneously valid prose punctuation and a structural marker. The instruction to ``write prose, not markdown'' passes over it because the em dash is already prose-legal.

Crucially, not all models exhibit this pattern. Our measurements show that Meta's Llama models produce zero em dashes across tens of thousands of words, while OpenAI and Anthropic models show elevated rates---with dramatic differences in how effectively suppression instructions reduce them. This variation is consistent with the em dash functioning not as a universal training artifact but as a \textbf{signature of the specific fine-tuning procedure applied}: different post-training methodologies amplify or suppress the latent tendency to different degrees.

\paragraph{Contributions.} This paper makes four contributions:

\begin{enumerate}[leftmargin=*]
\item A \textbf{five-step genealogy} for em dash production in LLM output, connecting training data composition, structural internalization, the dual-register status of the em dash, and post-training amplification into a single explanatory mechanism (Section~\ref{sec:genealogy}).

\item The \textbf{``last fingerprint'' framework}: a characterization of the em dash as the smallest unit of markdown-derived structural thinking that survives compression into prose, explaining why it persists under output constraints in models whose RLHF amplifies it (Section~\ref{sec:genealogy}).

\item \textbf{Empirical measurement} of em dash frequency and suppression resistance across twelve instruction-tuned LLMs from five providers ($\sim$240,000 words of generated prose), including a three-condition suppression gradient and a base-vs-instruct comparison, demonstrating that em dash behavior functions as a signature of fine-tuning methodology (Section~\ref{sec:empirical}).

\item A \textbf{documentation of two parallel discourses} about LLM writing behavior that have not been connected, and a demonstration that connecting them yields a unified explanation for both formatting-level and prose-level artifacts (Section~\ref{sec:discourses}).
\end{enumerate}
